\documentclass[11pt]{article}
% ------------------------------------------------
% Paquetes básicos
% ------------------------------------------------
\usepackage[a4paper,margin=1in]{geometry}
\usepackage{libertinus}
\usepackage{libertinust1math}
\usepackage[spanish]{babel}
\usepackage{xcolor}
\usepackage{hyperref}
\usepackage{titlesec}
\usepackage{tocloft}
\usepackage{amsmath, amssymb, amsthm}
\usepackage[most]{tcolorbox}
\usepackage{xifthen}
% ------------------------------------------------
% Colores
% ------------------------------------------------
\definecolor{indigo}{RGB}{40,60,120}
\definecolor{deepindigo}{RGB}{25,45,75}
% ------------------------------------------------
% Entornos
% ------------------------------------------------
\newtheoremstyle{ejstyle}
{1em}{1em}
{\normalfont}{}
{\bfseries\color{deepindigo}}{.}
{.5em}{}
{}
\theoremstyle{ejstyle}
\newtheorem{ejercicio}{Ejercicio}

\newenvironment{solucion}
{\begin{proof}[{\bfseries\color{deepindigo} Solución}]}
	{\end{proof}}

\newenvironment{obs}
{\par\smallskip\noindent{\bfseries\color{deepindigo} Observación.}\enspace\ignorespaces}
{\par\smallskip}
% ------------------------------------------------
% Comandos matemáticos
% ------------------------------------------------
\renewcommand{\P}{\mathbb{P}}
\newcommand{\E}{\mathbb{E}}
\newcommand{\R}{\mathbb{R}}
\newcommand{\Q}{\mathbb{Q}}
\newcommand{\Z}{\mathbb{Z}}
\newcommand{\C}{\mathbb{C}}

\newcommand{\mean}[2][]{%
	\ifthenelse{\isempty{#1}}%
	{\mathbb{E}\!\left[#2\right]}%
	{\mathbb{E}\!\left[#2\,\middle|\,\mathcal{F}_{#1}\right]}%
}
% ------------------------------------------------
\title{\Huge Entrega 3: Movimiento Browniano}
\author{Lorenzo Martinelli}
% ------------------------------------------------
\begin{document}
	
	\maketitle
	
	\begin{abstract}
		Tercer entrega de Cálculo Estocástico. Ejercicios 1, 3, 4, 7 y 9.
	\end{abstract}
	

	\begin{ejercicio}
		Sea $W(\cdot)$ un movimiento Browniano unidimensional. Mostrar que es una
		martingala.
	\end{ejercicio}
	
	\begin{solucion}
		Queremos ver que $\mean[s]{W(t)} = W(s)$ para todo $s \leq t$.
		\[
		\mean[s]{W(t)} = \mean[s]{W(s)} + \mean[s]{W(t)-W(s)}
		= W(s) + \mean{W(t)-W(s)} = W(s),
		\]
		donde usamos que $W(s)$ es $\mathcal{F}_s$-medible y que el incremento
		$W(t)-W(s)$ es independiente de $\mathcal{F}_s$ con $\mean{W(t)-W(s)}=0$.
	\end{solucion}
	
	\begin{ejercicio}
		Probar que si $W(\cdot)$ es un movimiento Browniano $n$-dimensional, entonces
		también lo son:
		\begin{enumerate}
			\item[\textit{a)}] $\mathbf{W}(t+s) - \mathbf{W}(s)$ para todo $s \geq 0$.
			\item[\textit{b)}] $c\,\mathbf{W}(t/c^2)$ para todo $c > 0$
			\quad (``rescale Browniano'').
		\end{enumerate}
	\end{ejercicio}
	
	\begin{solucion}
		Sea $X(t) = \mathbf{W}(t+s) - \mathbf{W}(s)$ con $s \geq 0$ fijo. Nuestra filtración va a ser $\mathcal{G}_t = \mathcal{F}_{t+s}.$
		\begin{enumerate}
			\item $X(0) = \mathbf{W}(s) - \mathbf{W}(s) = 0.$
			\item Para $r < t$,
			\[
			X(t) - X(r) = \mathbf{W}(t +s) - \mathbf{W}(t+s),
			\]
			que sabemos que es independiente de $\mathcal{F}_{r+s} = \mathcal{G}_s.$
			
			\item $X(t) - X(r) = \mathbf{W}(t+s) - \mathbf{W}(r+s)
			\sim \mathcal{N}\!\left(0,\,(t-r)I\right).$
			
			\item $t \mapsto \mathbf{W}(t+s)$ es continuo por
			serlo $\mathbf{W}(t)$.
		\end{enumerate}
		
		Veamos la otra. Sea $Y(t) = c\,\mathbf{W}(t/c^2)$, con $c > 0$ fijo. Con $\mathcal{G}_t = \mathcal{F}_{t/c^2}.$
		\begin{enumerate}
			\item $Y(0) = c\,\mathbf{W}(0) = \mathbf{0}.$
			\item Para $s < t$,
			\[
			Y(t) - Y(s) = c\bigl[\mathbf{W}(t/c^2) - \mathbf{W}(s/c^2)\bigr],
			\]
			independiente de $\mathcal{F}_{s/c^2} = \mathcal{G}_{s}$.
			
			\item Basta ver que 
			$$
			c\bigl[\mathbf{W}(t/c^2) - \mathbf{W}(s/c^2)\bigr] \sim c\,\mathcal{N}\!\left(0,\,\tfrac{t-s}{c^2}\,I\right)
					\sim \mathcal{N}\!\left(0,\,(t-s)I\right).
			$$
			
			\item Es continua pues $\mathbf{W}(t)$ lo es.
		\end{enumerate}
	\end{solucion}
	

	\begin{ejercicio}
		Sea $W(\cdot)$ un movimiento Browniano unidimensional. Mostrar que
		\[
		\lim_{k \to \infty} \frac{W(k)}{k} = 0 \qquad \text{casi seguramente.}
		\]
		\textit{Sugerencia:} Fijar $\varepsilon > 0$, definir el evento
		$A_k := \{|\frac{W(k)}{k}| \geq \varepsilon\}$ y aplicar Borel-Cantelli.
	\end{ejercicio}
	
	\begin{solucion}
		Fijamos $\varepsilon > 0$ y definimos $A_k$ como en la sugerencia.
		Como 
		$$ 
		\P(|\frac{W_k}{k}| \leq \varepsilon) = \P(|\frac{W_k}{\sqrt{k}}| \leq \varepsilon \sqrt{k}) = 2 \P(\frac{W_k}{\sqrt{k}} \leq \varepsilon \sqrt{k})
		$$
		Nos gustaría ver que esta probabilidad está acotada por una sucesión cuya serie converge pues Borel Cantelli nos daría el resultado. Notemos que $\frac{W_k}{\sqrt{k}} \sim \mathcal{N}(0,1)$. Luego, nuestra probabilidad está acotada por 
		$$
		\frac{1}{\sqrt{2\pi}}\int_{\sqrt{k}\varepsilon}^{\infty} \exp\{-\frac{x^2}{2}\}dx \leq \frac{1}{\sqrt{2\pi}}\int_{\sqrt{k}\varepsilon}^{\infty} \exp\{-\frac{\varepsilon\sqrt{k}x}{2}\}dx = \frac{2}{\sqrt{2\pi k}} \exp\{-\frac{\varepsilon k}{2}\}   
		$$
		Luego, cómo la serie de esta ultima expresión converge estamos.
	\end{solucion}
	
	\begin{ejercicio}
		Sea $X(t) := \int_0^t W(s)\,ds$, donde $W(\cdot)$ es un movimiento Browniano.
		Probar que
		\[
		\E\!\left[X^2(t)\right] = \frac{t^3}{3} \qquad \forall\, t > 0.
		\]
	\end{ejercicio}
	
	\begin{solucion}
		Por Fubini,
		\[
		\E\!\left[X^2(t)\right]
		= \E\!\left[\int_0^t W(s)\,ds \int_0^t W(u)\,du\right]
		= \int_0^t\!\int_0^t \E[W(s)W(u)]\,ds\,du.
		\]
		Cómo $\E[W(s)W(t)] = \min(s,t)$
		\[
		\E\!\left[X^2(t)\right]
		= \int_0^t\!\int_0^t \min(s,u)\,ds\,du
		= 2\int_0^t\!\int_0^u s\,ds\,du
		= 2\int_0^t \frac{u^2}{2}\,du
		= \int_0^t u^2\,du
		= \frac{t^3}{3}.
		\]
	\end{solucion}
	
	\begin{ejercicio}
		Sea $U(t) := e^{-t}W(e^{2t})$, donde $W(\cdot)$ es un movimiento Browniano
		unidimensional. Mostrar que
		\[
		\E[U(t)U(s)] = e^{-|t-s|} \qquad \text{para todo } -\infty < s,t < \infty.
		\]
	\end{ejercicio}
	
	\begin{solucion}
		Usando otra vez que $\E[W(s)W(t)] = \min(s,t)$,
		\[
		\E[U(t)U(s)]
		= e^{-t}e^{-s}\,\E\!\left[W(e^{2t})W(e^{2s})\right]
		= e^{-(t+s)}\min(e^{2t}, e^{2s}) = e^{-|t-s|}.
		\]
	\end{solucion}
	
\end{document}